\subsubsection{Feature Summary}
\label{subsubsec:feature_summary}
% ─────────────────────────────────────────────────────

\begin{table}[htbp]
\centering
\caption{Predictor set (26--35 features per fold, depending on the per-fold
PCA component count) shared by all tested baseline unimodal models, and
included in the newly proposed model. NTL is included only in the newly
proposed multimodal model.}
\label{tab:36_features}
\footnotesize
\begin{tabularx}{\columnwidth}{r l X}
\toprule
\textbf{\#} & \textbf{Feature} & \textbf{Derivation / Justification} \\
\midrule
\multicolumn{3}{l}{\textit{Temporal \& Calendar (12)}} \\
\addlinespace[2pt]
1  & \texttt{day\_of\_year}       & Ordinal seasonal position \\
2  & \texttt{week\_of\_year}      & Ordinal weekly position \\
3  & \texttt{quarter}             & Quarterly seasonal indicator \\
4  & \texttt{hour\_sin}           & Cyclical encoding ($r = -0.211$) \\
5  & \texttt{hour\_cos}           & Cyclical encoding ($r = -0.535$) \\
6  & \texttt{dow\_sin}            & Cyclical encoding ($r = 0.287$) \\
7  & \texttt{dow\_cos}            & Cyclical encoding ($r = -0.099$) \\
8  & \texttt{month\_sin}          & Cyclical encoding ($r = -0.112$) \\
9  & \texttt{month\_cos}          & Cyclical encoding ($r = 0.341$) \\
10 & \texttt{year}                & Trend (31.7\% of STL var.) \\
11 & \texttt{is\_public\_holiday} & $-12.1$\% load ($p < 10^{-300}$) \\
12 & \texttt{is\_school\_holiday} & $-2.5$\% load ($p \approx 10^{-140}$) \\
\midrule
\multicolumn{3}{l}{\textit{Autoregressive \& Cross-Correlation Lags (7)}} \\
\addlinespace[2pt]
13 & \texttt{load\_lag\_12h}        & ACF $r = 0.695$ \\
14 & \texttt{load\_lag\_24h}        & ACF $r = 0.842$ (diurnal) \\
15 & \texttt{load\_lag\_1w}         & ACF $r = 0.924$ (weekly) \\
16 & \texttt{load\_lag\_1y}         & ACF $r = 0.671$ (annual) \\
17 & \texttt{solar\_rad\_lag\_10h}  & CCF peak $-10$h ($r = -0.544$) \\
18 & \texttt{humidity\_lag\_12h}    & CCF peak $-12$h ($r = 0.411$) \\
19 & \texttt{temp\_lag\_107h}       & CCF peak $-107$h ($r = -0.399$) \\
\midrule
\multicolumn{3}{l}{\textit{PCA-Compressed Exogenous (7--16, refit per fold)}} \\
\addlinespace[2pt]
20\,--\,$19{+}k$ & \texttt{pca\_fold\_01\,...\,pca\_fold\_$k$} & $k \in [7, 16]$ leading PCs (${\geq}95$\% cumulative var.\ of 66 CBS+KNMI covariates, refit on each fold's training window) \\
\bottomrule
\end{tabularx}
\end{table}


\paragraph{Preprocessing at Training Time.}
Within each ROCV fold, all continuous features undergo IQR-based outlier removal (clipping at $Q_1 - 1.5 \times \text{IQR}$ and $Q_3 + 1.5 \times \text{IQR}$), followed by forward-filling as previously noted to handle any residual nulls. To optimize gradient-descent convergence, all continuous tabular variables are Min-Max normalized to a $[0, 1]$ range. \textit{Scaling parameters ($x_{\min}, x_{\max}$) are computed exclusively on the training fold} to prevent information leakage into validation and test sets. PCA is refitted within each fold on the 79 candidate exogenous CBS+KNMI covariates filtered for missingness over that fold's training window (retaining 66; up to 75 in the latest MLR folds, whose merged training$+$validation window extends furthest), and the minimal number of leading components capturing $\geq 95\%$ cumulative variance is kept (7--16 across folds), as described in \ref{sec:apx:second_appendix}.
